%% notes-tex/024-perturb-seq-pilot/sections/1-plan.tex
%% SOURCE: experiments/024-perturb-seq-costing/scripts/pilot_options.py
%%         (results/pilot_options.json; tables/t1-pilot-options.tex is emitted by it)
%% Every dollar figure composes cost_model.py's platforms over the Carver Center
%% rate card in uiuc_core_data.py. No price is entered here by hand.

\section{The Plan\secstatus{todo}}
\label{sec:plan}

Two runs, one per host, started in parallel by two people. Neither waits on the
other, and they answer different questions because the two organisms are at
different stages.

\notesec{\org{S.\ cerevisiae}: the library exists, so measure the readout} A
genome-wide single-guide CRISPR interference (\term{CRISPRi}) library is already
in hand, six guides per gene over roughly 6{,}000
genes~\citep{lianMultifunctionalGenomewideCRISPR2019}. What has never been
established is whether a cell's guide can be recovered from its RNA. Call that
probability \term{q}, the fraction of cells whose perturbation is read. One
channel on the core's 10x Genomics 5$'$ kit answers it (\cref{sec:guide}).

\notesec{\org{P.\ kudriavzevii}: no library is reachable, so measure the assay}
Transformation efficiency in this host is three to five orders of magnitude
below \org{S.\ cerevisiae}, which puts a pooled library-scale screen out of
reach until that moves. A first run therefore carries no guides at all, and what
it returns is whether the cell wall yields to in-droplet digestion, how many
messenger RNA molecules are recovered per cell, and what fraction of reads is
ribosomal. None of the three needs genetics.

%% GENERATED FILE -- do not hand-edit.
%% SOURCE: experiments/024-perturb-seq-costing/scripts/pilot_options.py
\begin{table}[htbp]\centering
\footnotesize
\caption[]{\textbf{The candidate first runs, priced end to end.} Usable cells are those expected to survive quality control, 55\% of a droplet channel's 20{,}000 and 25\% of a protocol run's 480{,}000. Recurring cost is reagents plus sequencing at the cheapest configuration the Carver Center sells that holds the reads, which at this scale is never the flow cell a genome-scale screen would use. The one-time column is the split-pool barcode plate set, which covers about 215 protocol runs. The route is the plate one for the eight-media row and 10x droplet for every other.}\label{tab:pilots}
\begin{tabular}{@{}>{\raggedright\arraybackslash}p{40mm}>{\raggedright\arraybackslash}p{23mm}rrr>{\raggedright\arraybackslash}p{44mm}@{}}
\toprule
run & host & usable cells & recurring & one-time & returns \\
\midrule
Guide capture, 1 channel & \org{S. cerevisiae} & 11,000 & \$4,493 & -- & q, the per-cell guide detection rate \\
Guide capture, low vs high copy & \org{S. cerevisiae} & 22,000 & \$8,576 & -- & q on two vector backbones \\
Unperturbed profile, 1 medium & \org{P. kudriavzevii} & 11,000 & \$3,810 & -- & wall digestion, UMIs per cell, rRNA fraction \\
Unperturbed profile, 4 media & \org{P. kudriavzevii} & 44,000 & \$12,550 & -- & the same three, plus a 4-condition response atlas \\
Plate route, 8 media in one run & \org{P. kudriavzevii} & 20,000 & \$4,990 & \$10,000 & the same three, across 8 media, plus the plate chemistry \\
\bottomrule
\end{tabular}
\end{table}


%% SOURCE: experiments/024-perturb-seq-costing/scripts/pilot_options.py
\tcfigfit[1.0]{figures/pilot_options.pdf}{%
  \textbf{The two first runs cost under \$9{,}000 together, and the cheaper way
  to add a growth medium is a plate rather than a second droplet channel.}
  \textbf{(a)} Recurring cost of each candidate, split into reagents (solid, by
  host) and sequencing (hatched). Sequencing is priced on the cheapest
  configuration the Carver Center sells that holds the reads, which at this
  scale is a MiSeq run or a split NovaSeq lane rather than the large flow cell a
  genome-scale screen fills. The plate route carries a \$10{,}000 one-time
  purchase of barcoded oligo plates, which covers about 215 protocol runs.
  \textbf{(b)} Recurring cost against the number of growth media, at one droplet
  channel's worth of cells per medium. An unmodified droplet channel carries
  nothing that says which medium a cell came from, so each medium buys its own
  channel, at \$2{,}650 for the first and \$2{,}240 for each of the next six.
  Both plate routes read a
  well index out of every cell already, so media ride together in one run and
  the steps are sublibrary and run rounding rather than per-medium cost.
  \textbf{(c)} What one run buys against what a screen needs, on usable cells.
  \textbf{(d)} Why $q$ is the quantity worth measuring first. At one guide per
  cell it scales the yield linearly; at $k$ guides per cell the usable fraction
  is $q^{k}$, so the two published microbial values, 0.21 in
  \org{E.\ coli}~\citep{brandnerPooledSinglecellCRISPRa2025} and more than 0.71
  in \org{S.\ cerevisiae}~\citep{nadal-ribellesSinglecellResolvedGenotypephenotype2025},
  are the difference between multiplexing being expensive and being impossible.
  Generated by \file{experiments/024-perturb-seq-costing/scripts/pilot\_options.py}.}{fig:pilots}

\notesec{What the prices do and do not include} Reagents and sequencing at
on-campus rates, effective 2026-08-01. Not included: strain construction, the
media themselves, staff time, and the cell counter discussed in
\cref{sec:counting}. The sequencing line is the one most likely to move, because
a lane is an indivisible purchase and a run that half-fills one pays for the
whole thing.

\begin{keybox}
\textbf{Start both. They cost \$8{,}303 together and answer different
questions.} In \org{S.\ cerevisiae}, one 5$'$ channel on the existing CRISPRi
library measures $q$, the fraction of cells whose guide is readable, at
\$4{,}493. In \org{P.\ kudriavzevii}, one channel with no library measures wall
digestion, messenger RNA per cell, and ribosomal fraction, at \$3{,}810. Neither
depends on the other, one person can own each, and the two results together
decide every design question after them.
\end{keybox}
